\section{Backend изпълнителни модели и native client библиотеки}
\label{sec:backend_scenarios}

Настоящата глава изследва не само как една и съща C++ абстракция се материализира в различни backend-и, но и как се изпълнява през конкретните native client библиотеки. Анализът има две цели. Първо, да покаже, че библиотеката не е ограничена до чисто SQL интерпретация. Второ, да изясни къде асинхронната интеграция може да следва native non-blocking модел и къде честният архитектурен избор остава thread-pool offload върху синхронен драйвер.

След теоретичната рамка от Глава~\ref{sec:theoretical_background} и архитектурното обяснение на async слоя от Глава~\ref{sec:async_execution_architecture}, тук вече не се пита дали различните модели за съхранение са еквивалентни. Вместо това се пита дали една и съща логическа ORM форма може да бъде адаптирана контролирано към тях, без библиотеката да обещава execution semantics, които даден backend не може естествено да изпълнява. Тъкмо в това се проявява стойността на capability-базирания connector contract и на разграничението между универсален async интерфейс и backend-специфична реализация.

\subsection{Методика на анализа}

Всеки сценарий следва една и съща последователност: теоретична опора, логическо картографиране от C++ модела, реализация в \code{connector\_trait<DB>}, ограничения и верификация чрез тестове. Това прави сравнението между backend-ите възможно. Не се търси твърдение, че всички системи са еквивалентни, а че една и съща логическа ORM абстракция може да бъде адаптирана контролирано към различни модели за съхранение.

\begin{table}[H]
\centering
\small
\caption{Еднаква аналитична рамка за всеки backend сценарий}
\label{tab:backend_analysis_frame}
\begin{tabularx}{\textwidth}{L{3.3cm}Y}
\toprule
\textbf{Критерий} & \textbf{Какво се оценява} \\
\midrule
Теоретична опора & към кой модел за съхранение от Глава~\ref{sec:theoretical_background} принадлежи backend-ът \\
Логическо картографиране & как \code{property}, \code{field}, \code{Rule} и \code{table\_name\_trait} се интерпретират в конкретния backend \\
Capability профил & кои ORM операции са изрично допустими и кои трябва да бъдат отхвърлени \\
Execution поведение & как се рендират заявки, параметри и резултатни проекции \\
Тестова опора & кои unit и integration тестове доказват contract-а в хранилището \\
\bottomrule
\end{tabularx}
\end{table}

Тази аналитична рамка е важна, защото позволява сравнението да остане честно. SQL backend-ите не се привилегироват само защото имат по-богат синтаксис, а NoSQL backend-ите не се оценяват като ``непълни'' само защото отказват операции, които са неподходящи за техния модел на съхранение.

\begin{figure}[H]
\centering
\begin{tikzpicture}[node distance=1.35cm and 1.2cm]
    \node[ormaccent] (ir) {\textbf{Споделен query IR}\\fields, predicates, placeholders, projections};
    \node[ormblock, below left=of ir] (sql) {SQL материализация\\SQLite/PostgreSQL/MySQL};
    \node[ormblock, below=of ir] (doc) {Документна материализация\\MongoDB};
    \node[ormblock, below right=of ir] (graph) {Graph materialization\\Neo4j};
    \node[ormblock, below=of doc, xshift=-3.2cm] (redis) {Key-based materialization\\Redis};
    \node[ormblock, below=of doc, xshift=3.2cm] (cass) {Partition-aware materialization\\Cassandra};

    \draw[ormline] (ir) -- (sql);
    \draw[ormline] (ir) -- (doc);
    \draw[ormline] (ir) -- (graph);
    \draw[ormline] (ir) -- (redis);
    \draw[ormline] (ir) -- (cass);
\end{tikzpicture}
\caption{Един споделен query IR и множество backend-specific materialization пътища}
\label{fig:backend_materialisation_paths}
\end{figure}

Фигура~\ref{fig:backend_materialisation_paths} показва ключовото твърдение на главата: не съществуват множество ORM библиотеки в рамките на хранилището, а една логическа библиотека с различни материализационни пътища. Това е силният аргумент в полза на connector architecture-а --- backend разнообразието е овладяно чрез общ IR и explicit contract, а не чрез копиране на API-та.

\subsection{Релационен базов сценарий: SQLite}

Този сценарий стъпва върху релационния модел, разгледан в Глава~\ref{sec:theoretical_background}. \code{SQLiteDB} е най-естественият референтен backend за библиотеката, защото съчетава реална SQL семантика, малък operational контекст и пряка интеграция чрез \code{sqlite3} API. В него \code{property} се материализира като колона, \code{table\_name\_trait} --- като име на таблица, а \code{Rule} дърветата --- като \code{WHERE} изрази.

Конекторът декларира \code{supports\_joins}, \code{supports\_transactions} и \code{supports\_aggregation}, което го прави най-близък до класическото ORM очакване. Именно върху него се валидират подготвени заявки, bind логика, \code{find\_one()} сценарии и result hydration към typed tuples. Интеграционните тестове в \path{tests/integration/test_sqlite.cpp} показват не само коректност на \code{SELECT}, \code{INSERT}, \code{UPDATE} и \code{DELETE}, но и поведение при transaction-capable connector и join-capable profile.

От академична гледна точка SQLite играе ролята на релационен baseline. Ако дадена ORM форма е проблематична още тук, тя трудно би могла да бъде обобщена към по-екзотични backend-и. Затова настоящата работа използва SQLite като отправна точка, а не просто като един от многото примери.

\subsection{Релационна преносимост: PostgreSQL и MySQL}

Теоретичната основа и тук е релационният модел, но архитектурният интерес е различен: доколко един и същ query IR остава валиден при различни SQL диалекти и различни правила за bind на параметри. \code{PostgreSQLDB} и \code{PostgreSQLLiveDB} поддържат \code{supports\_joins}, \code{supports\_transactions} и \code{supports\_aggregation}, а тестовете в \path{tests/unit/test_postgresql_connector.cpp} и \path{tests/integration/test_postgresql_live.cpp} показват реално изпълнение срещу жива база данни.

\code{MySQLDB} и \code{MySQLLiveDB} също декларират \code{supports\_joins} и \code{supports\_transactions}, но тук е важно различието в bind модела. Докато PostgreSQL може естествено да адресира параметри чрез \code{$1}, \code{$2}, \ldots{}, MySQL разчита на позиционни \code{?} placeholders. Това не нарушава общия IR, но показва, че logical portability не означава byte-for-byte идентична backend материализация.

Тъкмо тези SQL family сценарии защитават по-силна теза от ``работи със SQL'': библиотеката демонстрира, че query IR-ът е достатъчно общ за повече от един диалект, а в същото време е достатъчно честен, за да остави нисконивовите различия на connector слоя.

\subsection{Документен сценарий: MongoDB}

Този сценарий стъпва върху документния модел от теоретичната глава. При MongoDB \code{property} се интерпретира като поле в документ, а query predicate дърветата се превеждат към BSON-подобни филтри. \code{select} изразите се материализират като projection, а таблицата от ORM гледна точка се превръща в collection. Точно тук се вижда практическата стойност на storage-agnostic query IR: потребителят не описва \code{\$eq}, \code{\$and} и projection директно, но конекторът може да ги изведе от същата логическа форма.

Важна подробност е, че \code{connector\_trait<MongoDB>} не декларира SQL-style capability tags. Това не е пропуск, а коректно описание на connector contract-а. Няма претенция, че Mongo backend-ът е join-equivalent на релационните конектори. Вместо това той предлага документно-естествено тълкуване на query IR-а. Unit и integration тестовете в \path{tests/unit/test_mongodb_connector.cpp} и \path{tests/integration/test_mongodb_live.cpp} валидират тази адаптация срещу реален driver слой.

Архитектурната стойност на този сценарий е, че доказва едно важно ограничение: една библиотека може да бъде multi-backend, без да твърди, че всички backend-и поддържат едни и същи операции. Именно този тип ``контролирана частичност'' прави дизайна научно защитим.

\subsection{Key-value сценарий: Redis}

Този сценарий стъпва върху key-value теорията. При Redis един и същ entity модел може да бъде материализиран по два основни начина: като \code{SET}/\code{GET} при еднополеви структури или като \code{HSET}/\code{HGETALL} при многополеви структури. \code{table\_name()} участва в генерирането на key prefix, а достъпът е концентриран около primary-key equality.

Тук ясно се вижда, че споделеният query IR не означава идентична изразителност. \code{RedisDB} умишлено не декларира \code{supports\_joins}, \code{supports\_transactions} или \code{supports\_aggregation}. Това ограничение е част от коректния contract, а не слабост на архитектурата. Unit и live тестовете в \path{tests/unit/test_redis_connector.cpp} и \path{tests/integration/test_redis_live.cpp} доказват, че библиотеката правилно материализира key-based достъп, без да обещава неподдържани операции.

Този сценарий е особено ценен за дипломната работа, защото принуждава архитектурата да прави ясно разграничение между \emph{логически модел на данните} и \emph{допустим модел на достъп}. Това разграничение е лесно да се забрави при SQL-first библиотеките, но тук е изведено като централен принцип.

\subsection{Графов сценарий: Neo4j}

Сценарият за Neo4j стъпва върху графовия модел. Entity типът се материализира като label и набор от свойства, а query IR се превежда към \code{MATCH}, \code{RETURN} и \code{WHERE} фрагменти. Това е важен архитектурен тест, защото graph query семантиката е най-отдалечена от класическия SQL модел. Въпреки това логическите елементи --- полета, предикати, филтриране и резултатни проекции --- остават разпознаваеми.

\code{Neo4jDB} и \code{Neo4jLiveDB} декларират \code{supports\_transactions}, но не и join/aggregation capabilities на SQL слоя. Това е смислено: графовият traversal не е същото като SQL join, макар на логическо ниво и двете да свързват данни. Unit и integration тестовете в \path{tests/unit/test_neo4j_connector.cpp} и \path{tests/integration/test_neo4j_live.cpp} показват, че abstraction-ът може да адаптира един и същ логически модел и към graph backend, ако конекторът поеме отговорността за правилната материализация.

От гледна точка на тезата това е една от най-силните демонстрации, че query IR-ът е семантичен, а не синтактичен слой. Ако IR-ът беше просто недо-SQL, Neo4j сценарият би се разпаднал. Фактът, че е възможна работеща графова материализация, подкрепя централната архитектурна претенция.

\subsection{Wide-column сценарий: Cassandra}

Този сценарий стъпва върху wide-column теорията. Макар Cassandra да предлага CQL синтаксис, наподобяващ SQL, действителният модел на достъп е подчинен на partition ключовете и на разпределената организация на данните. Затова библиотеката не може да третира Cassandra като обикновен SQL backend. Наличието на \code{WHERE} клауза само по себе си не е достатъчно; тя трябва да бъде съвместима с partition логиката.

\code{CassandraDB} и \code{CassandraLiveDB} декларират \code{supports\_transactions}, но не и \code{supports\_joins}. Това е много показателно. От една страна, backend-ът продължава да е част от същата ORM библиотека. От друга страна, contract-ът категорично отказва операции, които биха били подвеждащи или скъпи в wide-column среда. Unit и live тестовете в \path{tests/unit/test_cassandra_connector.cpp} и \path{tests/integration/test_cassandra_live.cpp} дават конкретна верификация на този по-строг operational profile.

Именно поради това Cassandra е най-добрият пример, че разширяемостта на библиотеката не е базирана на фалшиво уеднаквяване, а на формално описани граници на допустимите операции.

\begin{figure}[H]
\centering
\begin{tikzpicture}[node distance=1.2cm and 0.95cm]
    \node[ormaccent, minimum width=3.4cm] (title) {Capability спектър};
    \node[ormblock, below left=1.5cm and 2.3cm of title, minimum width=2.8cm] (sqlite) {SQLite\\join + txn + agg};
    \node[ormblock, below=1.5cm of title, minimum width=3.1cm] (pgsql) {PostgreSQL/MySQL\\join + txn + agg/txn};
    \node[ormblock, below right=1.5cm and 2.3cm of title, minimum width=2.8cm] (mongo) {MongoDB\\без SQL capability tags};
    \node[ormblock, below=of sqlite, minimum width=2.8cm] (redis) {Redis\\key-based only};
    \node[ormblock, below=of pgsql, minimum width=2.8cm] (neo4j) {Neo4j\\txn only};
    \node[ormblock, below=of mongo, minimum width=2.8cm] (cass) {Cassandra\\txn only};

    \draw[ormline] (title) -- (sqlite);
    \draw[ormline] (title) -- (pgsql);
    \draw[ormline] (title) -- (mongo);
    \draw[ormline] (sqlite) -- (redis);
    \draw[ormline] (pgsql) -- (neo4j);
    \draw[ormline] (mongo) -- (cass);
\end{tikzpicture}
\caption{Сравнителен capability профил на backend-ите в хранилището}
\label{fig:backend_capability_spectrum}
\end{figure}

Фигура~\ref{fig:backend_capability_spectrum} синтезира едно от най-важните наблюдения в главата: библиотеката не налага еднакъв capability профил на всички backend-и. Точно обратното --- различията се използват като first-class архитектурна информация.

\subsection{Синтезирана capability матрица}

За да остане сравнението не само качествено, но и таблично проследимо, Фигура~\ref{fig:backend_capability_spectrum} може да бъде сведена до компактна capability матрица. Тя не бива да се чете като универсална истина за самите системи за бази данни, а като описание на конкретната архитектурна позиция на библиотеката спрямо тях в разглежданата ревизия на хранилището.

\begin{table}[H]
\centering
\small
\caption{Capability профили на основните backend-и според наличните тестове и connector specializations}
\label{tab:backend_capability_matrix}
\begin{tabularx}{\textwidth}{L{2.0cm}L{1.45cm}L{1.55cm}L{1.5cm}L{1.35cm}L{1.55cm}Y}
\toprule
\textbf{Backend} & \textbf{Joins} & \shortstack{\textbf{Транзак-}\\\textbf{ции}} & \shortstack{\textbf{Агрега-}\\\textbf{ция}} & \textbf{Upsert} & \shortstack{\textbf{Bulk}\\\textbf{insert}} & \textbf{Кратък коментар} \\
\midrule
SQLite & налична & налична & налична & \shortstack{ограни-\\чена} & \shortstack{ограни-\\чена} & най-близкият пълен релационен reference path в хранилището \\
PostgreSQL & налична & налична & налична & \shortstack{специ-\\фична} & \shortstack{въз-\\можна} & силен релационен contract с live execution evidence \\
MySQL & налична & налична & налична & \shortstack{специ-\\фична} & \shortstack{въз-\\можна} & operational path е наличен, но bind и dialect детайлите остават connector-специфични \\
MongoDB & липсва & липсва & липсва & аналог & \shortstack{не е\\централна} & силен document model без SQL-style capability профил \\
Redis & липсва & липсва & липсва & аналог & \shortstack{ограни-\\чена} & key-value semantics с минималистичен, но честен contract профил \\
Neo4j & аналог & налична & липсва & \shortstack{не е\\централна} & \shortstack{не е\\централна} & graph backend с transaction focus, а не със SQL join/aggregation слой \\
Cassandra & липсва & налична & липсва & \shortstack{не е\\централна} & \shortstack{въз-\\можна} & wide-column backend с ограничен, но формално дефиниран capability набор \\
\bottomrule
\end{tabularx}
\end{table}

Таблица~\ref{tab:backend_capability_matrix} позволява по-прецизно да се види къде библиотеката е най-близо до класическо ORM поведение и къде съзнателно работи с частичен contract. Релационните backend-и концентрират по-богат capability профил, докато документните, графовите и key-value вариантите се оценяват през други критерии. Именно тази контролирана нееднородност прави multi-backend дизайна инженерно честен.

\subsection{Съпоставка между test evidence и backend contract}

\begin{table}[H]
\centering
\small
\caption{Backend-и, capability профил и налична тестова опора}
\label{tab:backend_evidence}
\begin{tabularx}{\textwidth}{L{2.1cm}Q{2.4cm}L{2.8cm}Y}
\toprule
\textbf{Backend} & \textbf{Capability профил} & \textbf{Unit/contract тест} & \textbf{Integration/live тест} \\
\midrule
SQLite & joins, transactions, aggregation & capability asserts в \path{tests/integration/test_sqlite.cpp} & \path{tests/integration/test_sqlite.cpp} \\
PostgreSQL & joins, transactions, aggregation & \path{tests/unit/test_postgresql_connector.cpp} & \path{tests/integration/test_postgresql_live.cpp} \\
MySQL & joins, transactions, aggregation & \path{tests/unit/test_mysql_connector.cpp} & \path{tests/integration/test_mysql_live.cpp} \\
MongoDB & без SQL capability tags & \path{tests/unit/test_mongodb_connector.cpp} & \path{tests/integration/test_mongodb_live.cpp} \\
Redis & без joins/transactions/aggregation & \path{tests/unit/test_redis_connector.cpp} & \path{tests/integration/test_redis_live.cpp} \\
Neo4j & transactions & \path{tests/unit/test_neo4j_connector.cpp} & \path{tests/integration/test_neo4j_live.cpp} \\
Cassandra & transactions & \path{tests/unit/test_cassandra_connector.cpp} & \path{tests/integration/test_cassandra_live.cpp} \\
\bottomrule
\end{tabularx}
\end{table}

Таблица~\ref{tab:backend_evidence} е важна, защото превежда сравнението от равнището на narrative analysis към равнището на проверима repository evidence. За всеки backend има поне един ясен contract-level или live-level корелат. Това намалява риска главата да се превърне в декларативен преглед без техническа опора.

\subsection{Сравнителен синтез}

Съпоставката между backend-ите показва, че библиотеката успява да запази единен логически модел на равнището на entity типове, полета, placeholder-и и query IR. Различията се появяват на равнището на capability и на физическата организация на данните. Това е и най-важният извод на настоящата глава: успешната multi-backend ORM архитектура не премахва различията между моделите за съхранение, а ги прави контролируеми чрез формален connector contract.

\begin{table}[H]
\centering
\small
\caption{Съпоставка на backend материализацията}
\label{tab:backend_materialisation}
\begin{tabularx}{\textwidth}{L{2.0cm}L{2.5cm}L{2.8cm}Y}
\toprule
\textbf{Backend} & \textbf{Основна структура} & \textbf{Логическа интерпретация} & \textbf{Показателно ограничение} \\
\midrule
SQLite & таблица & пълна SQL ORM интерпретация & стандартна релационна дисциплина и богата query изразителност \\
PostgreSQL/MySQL & таблица & преносим SQL IR & различен bind модел и диалектни особености \\
MongoDB & документ & filter/projection интерпретация & частична еквивалентност спрямо SQL операции \\
Redis & ключ/стойност, hash & key-based entity access & lookup предимно по primary key и липса на join semantics \\
Neo4j & възел/свойства & graph-pattern интерпретация & graph-специфична query семантика вместо SQL join слой \\
Cassandra & wide-column row & CQL с ключови ограничения & partition-driven достъп и ограничен predicate profile \\
\bottomrule
\end{tabularx}
\end{table}

В резултат multi-backend design-ът на библиотеката може да бъде защитен с по-прецизна формулировка. Системата не обещава ``универсална база данни през един API'', а \emph{унифициран логически слой с експлицитни backend граници}. Именно тази формулировка е едновременно инженерно честна и научно обоснована.
